% Post-checkpoint-13 fragment; full mathematical and primary-source audits passed.
\subsection{Second cohomology and central lifting obstructions}

\begin{theorem}[Proved: a uniform cross-characteristic second-cohomology bound]
\label{D:surface-tame-cross-characteristic-H2}
Let $S$ be a normal proper connected surface over an algebraically
closed field $k$ of characteristic $p>0$, let $T\subsetneq S$ be
closed, and put $U=S-T$ and $Q=Q(S,T)$.  Thus $Q$ requires tameness
only at the original codimension-one boundary valuations.
There is a constant $c=c(S,T)$ such that, for every finite field
$\Lambda$ of characteristic $\ell\ne p$ and every finite-dimensional
continuous $\Lambda[Q]$-module $M$,
\[
          \dim_\Lambda H^2(Q,M)\le c\dim_\Lambda M.       \tag{D.H2.1}
\]
For an open subgroup $H\subset Q$ and a finite-dimensional continuous
$\Lambda[H]$-module $M$, the corresponding bound is
\[
          \dim_\Lambda H^2(H,M)
                \le c[Q:H]\dim_\Lambda M.               \tag{D.H2.2}
\]
All group cohomology is continuous.
\end{theorem}

\begin{proof}
Write $G=\pi_1(U)$ and $N=\ker(G\to Q)$.  By the definition of
this quotient, $N$ is the closed subgroup generated by the
$G$-conjugates of the images of the wild inertia groups at the
original boundary divisors.  Those groups are pro-$p$.
The $N$-action on $M$ is trivial.  A continuous homomorphism
from $N$ to the additive $\ell$-group $M$ vanishes on every
one of the generating pro-$p$ groups, and hence vanishes.
Consequently
\[
                        H^1(N,M)=0.
\]
The degree-two inflation map therefore gives an injection
\[
                H^2(Q,M)\lhook\joinrel\longrightarrow H^2(G,M).
                                                        \tag{D.H2.3}
\]
This is the inflation criterion in
\citep[II, Proposition~1.34 and Remark~1.35, pp.~71--72;
continuous cochains in II, Proposition~4.2, p.~87]{MilneCFT2020D}.
One can also check the injection directly: a continuous
$M$-extension of $Q$ whose pullback splits over $G$ has a
splitting $G\to E$.  Its restriction to $N$ is a homomorphism
$N\to M$, hence zero, so the splitting descends to $Q$.
This argument also explains why arbitrary prime-to-$p$
finite abelian coefficients will be allowed below.

Let $\mathcal M$ be the locally constant sheaf on $U$
associated with $M$.  There is a second injection
\[
              H^2(G,M)\lhook\joinrel\longrightarrow
                        H^2_{\mathrm{et}}(U,\mathcal M).
                                                        \tag{D.H2.4}
\]
For completeness, use the Cartan--Leray spectral sequence
for the system of all connected finite Galois covers of
$U$, as in \citep[Theorem~14.9, Remark~14.10 and
Proposition~14.13, pp.~96--98]{MilneLEC2013H2D}.
Its geometric degree-zero term is $M$.  Its degree-one
term is zero: after trivializing the coefficient sheaf,
every degree-one class is a finite torsor and dies on a
further finite cover.  In the first-quadrant spectral
sequence, the only possible incoming differential to
$(2,0)$ is therefore zero.  There are no outgoing
differentials from this term.  Its limiting term embeds
as the last filtration piece of total degree two,
which proves (D.H2.4).  We make no assertion that all
higher cohomology is group cohomology.

Choose a fixed smooth projective resolution
$\pi:Y\to S$ such that the reduced inverse image of $T$
and the exceptional curves have simple normal crossings.
Write $V=\pi^{-1}(U)$ and $h:V\to U$.  The map $h$ is
proper and birational, and it is an isomorphism outside
a finite set $\Sigma\subset U$.  The proper birational
map to the normal surface has connected geometric fibres.
Proper base change therefore gives
\[
             h_*h^*\mathcal M=\mathcal M,
\qquad
 (R^1h_*h^*\mathcal M)_{\bar x}
     =M_{\bar x}\otimes_\Lambda H^1(Y_{\bar x},\Lambda).
                                                        \tag{D.H2.5}
\]
The second identity can equivalently be checked after
etale-locally trivializing $\mathcal M$.  We use proper
base change in the precise form
\citep[Theorems~17.7 and~17.10, pp.~116--117]
{MilneLEC2013H2D}.
The sheaf in degree one is supported on $\Sigma$.
Nilpotents on the geometric fibres do not affect their
etale cohomology.  For their reduced simple-normal-crossing
curves, normalization and the node exact sequence show
\[
 \dim_\Lambda H^1(Y_{\bar x},\Lambda)
     =2\sum_jg(E_{x,j})+b_1(\Gamma_x),                    \tag{D.H2.6}
\]
where $E_{x,j}$ are the smooth components and $\Gamma_x$
is their dual graph; a zero-dimensional fibre contributes
zero.  In particular the sum $b$ of these dimensions over
$x\in\Sigma$ is independent of $\Lambda$.

The Leray spectral sequence for $h$
\citep[Theorem~12.7, pp.~82--83]{MilneLEC2013H2D}
now implies
\[
 \dim_\Lambda H^2(U,\mathcal M)
       \le \dim_\Lambda H^2(V,h^*\mathcal M)
                                +b\dim_\Lambda M.       \tag{D.H2.7}
\]
Here the direction is essential.  The term
$E_2^{2,0}=H^2(U,\mathcal M)$ has only the possible
incoming differential from $E_2^{0,1}$; its quotient
by that image is $E_\infty^{2,0}$, a subspace of
$H^2(V,h^*\mathcal M)$.  Thus no degeneration
assertion, and no vanishing of $R^2h_*$, is needed.

By Theorem~\ref{D:surface-uniform-conductor}, the
logarithmic conductors of $h^*\mathcal M$ along the
fixed boundary $Y-V$ have a bound independent of
$\ell$, the coefficient field, and the rank.
The theorem applies to the finite quotient cover
through which the representation factors; a bound
on its ramification filtration bounds every such
module.  At the strict transforms of the original
divisors the ramification is tame, and the remaining
components have the proved uniform exceptional bound.
The passage from these generic logarithmic bounds to the
curvewise bound required by the Betti theorem is exactly
Proposition~\ref{D:surface-curvewise-conductor}, using
\citep[Proposition~5.8(3)]{C_HuTeyssierRam2025}.
Hu--Teyssier's finite-coefficient theorem therefore
provides $b'$ depending only on this fixed pair with
\[
                   \dim_\Lambda H^2(V,h^*\mathcal M)
                            \le b'\dim_\Lambda M;
\]
see \citep[Theorem~7.34 and its proof, p.~41]
{C_HuTeyssierBetti2025}.  This is its $j=2$ case on
a smooth proper compactification; if the boundary
is empty the same proper smooth Betti bound applies.
Combining the two injections and (D.H2.7) proves
(D.H2.1), with $c=b+b'$.  Finally Shapiro's lemma
identifies $H^2(H,M)$ with
$H^2(Q,\operatorname{Ind}_H^Q M)$, whose coefficient
rank is $[Q:H]\dim M$, proving (D.H2.2).
\end{proof}

\begin{corollary}[Proved: polynomially many central obstruction classes]
\label{D:surface-tame-central-obstructions}
In the preceding theorem we may increase $c$ to an integer so
that, for every finite continuous abelian $Q$-module $A$ of
order prime to $p$,
\[
                         |H^2(Q,A)|\le |A|^c.            \tag{D.H2.8}
\]
In particular, for $q=p^f$ and $d=\gcd(r,q-1)$, the central
obstructions to lifting homomorphisms
$Q\to\operatorname{PSL}_r(\mathbf F_q)$ to
$\operatorname{SL}_r(\mathbf F_q)$ take at most $d^c$
values.  The obstructions for
$Q\to\operatorname{PGL}_r(\mathbf F_q)$ and its
$\operatorname{GL}_r(\mathbf F_q)$ lift take at most
$(q-1)^c$ values.

For every fixed integer $B$ there is a single open normal
subgroup $H_B\subset Q$ such that restriction to $H_B$
kills every class in $H^2(Q,\mathbf Z/d\mathbf Z)$
for every $1\le d\le B$ prime to $p$.
\end{corollary}

\begin{proof}
A finite $Q$-module of order prime to $p$ admits a
composition series whose factors are finite-dimensional
$\mathbf F_\ell[Q]$-modules for primes $\ell\ne p$.
For a short exact sequence $0\to A'\to A\to A''\to0$,
the cohomology exact sequence gives
\[
                 |H^2(Q,A)|
                   \le |H^2(Q,A')|\,|H^2(Q,A'')|.
\]
Apply (D.H2.1) to the factors and multiply these
inequalities.  Their cardinalities multiply to $|A|$,
so (D.H2.8) follows.

The two central kernels are respectively the cyclic
scalar groups of order $d$ and $q-1$, with trivial
$Q$-action.  Pulling back the corresponding central
extension defines its continuous degree-two
obstruction class, which vanishes exactly when a
continuous lift exists.  This is the cocycle
classification of extensions
\citep[II, Example~1.18(b), pp.~65--66, and
Proposition~4.2, p.~87]{MilneCFT2020D}.
Thus (D.H2.8) gives the two stated cardinalities.
These count obstruction classes, not representations
within one obstruction class.

For fixed $B$, there are finitely many relevant
coefficient groups and finitely many classes for
each by (D.H2.8).  Every class is represented by a
normalized continuous cocycle factoring through
one finite quotient of $Q$, by the continuous-cochain
statement of the same Proposition~4.2.  Its restriction
to the kernel of that quotient is zero.  Intersect
these finitely many open normal kernels to obtain
$H_B$.
\end{proof}

\begin{remark}[The remaining projective-representation problem]
For unbounded central order, the number of obstruction classes
is now controlled, but a uniform count inside one nontrivial
class is a separate problem.  Applying the linear theorem to
an endomorphism algebra replaces rank $r$ by $r^2$ and yields
an exponent of order $r^4$.  A rank-$r$ twisted unit-bundle
presentation with constants uniform in the central class and
its order has not been supplied by this argument.  No bound
on arbitrary nonabelian crowns or full finite generation is
inferred from (D.H2.1)--(D.H2.8).
\end{remark}
